\documentclass[11pt,a4paper]{article}

% ---- Packages ----
\usepackage[margin=2.5cm]{geometry}
\usepackage[utf8]{inputenc}
\usepackage[T1]{fontenc}
\usepackage{lmodern}
\usepackage{amsmath,amssymb,amsthm}
\usepackage{graphicx}
\usepackage[dvipsnames]{xcolor}
\usepackage{booktabs}
\usepackage{tabularx}
\usepackage{longtable}
\usepackage{multirow}
\usepackage{enumitem}
\usepackage{caption}
\usepackage[numbers,sort&compress]{natbib}
\usepackage{titlesec}
\usepackage{fancyhdr}
\usepackage[hidelinks,colorlinks=true,linkcolor=blue!60!black,citecolor=blue!60!black,urlcolor=blue!60!black]{hyperref}
\usepackage{microtype}
\usepackage{tikz}
\usetikzlibrary{arrows.meta,positioning}
\usepackage{pgfplots}
\pgfplotsset{compat=1.18}

% ---- Page style ----
\pagestyle{fancy}
\fancyhf{}
\renewcommand{\headrulewidth}{0.4pt}
\fancyhead[L]{\small Supplementary Information}
\fancyhead[R]{\small Yang (2026)}
\fancyfoot[C]{\thepage}

% ---- Section formatting ----
\titleformat{\section}{\large\bfseries}{S\thesection}{0.5em}{}
\titleformat{\subsection}{\normalsize\bfseries}{S\thesection.\thesubsection}{0.5em}{}

% ---- Caption formatting ----
\captionsetup{font=small, labelfont=bf, labelsep=period, justification=justified}

% ---- Paragraph style ----
\setlength{\parindent}{0pt}
\setlength{\parskip}{6pt}

% ---- Theorem environments ----
\newtheorem{theorem}{Theorem}
\newtheorem{definition}[theorem]{Definition}
\newtheorem{proposition}[theorem]{Proposition}
\newtheorem{lemma}[theorem]{Lemma}
\newtheorem{remark}[theorem]{Remark}

% ---- Custom commands ----
\newcommand{\R}{\mathbb{R}}
\newcommand{\Z}{\mathbb{Z}}
\newcommand{\Sclass}{\mathfrak{S}}
\newcommand{\norm}[1]{\left\|#1\right\|}
\newcommand{\Spec}{\mathcal{S}}
\newcommand{\Agent}{\mathcal{A}}
\newcommand{\Domain}{\Omega}
\newcommand{\cmark}{\checkmark}

\begin{document}

\begin{center}
{\Large\bfseries Supplementary Information}\\[12pt]
{\large A Judge Agent Closes the Reliability Gap in AI-Generated Scientific Simulation}\\[8pt]
{\normalsize Chengshuai Yang}\\[4pt]
{\small NextGen PlatformAI C Corp, USA}
\end{center}

\vspace{12pt}
\tableofcontents
\newpage

% =====================================================================
\section{Error Bound Proof within the Simulability Class $\Sclass$}
\label{sec:s1}
% =====================================================================

We restate and prove the main formal guarantee, which applies to problems in the simulability class $\Sclass$ (Definition~1, main text).

\begin{proposition}[Automated Bounded-Error Realizability within $\Sclass$]
\label{prop:bound}
Under the following assumptions, which correspond to the four conditions of $\Sclass$ (Definition~1, main text):
\begin{enumerate}[nosep,label=(A\arabic*)]
\item \textbf{Finite specifiability (S1):} $P$ admits a finite representation $\Spec = (\Domain, \mathcal{E}, \mathcal{B}, \mathcal{I}, \mathcal{O}, \varepsilon)$, realized as a valid \texttt{spec.md} file. The Plan Agent generates this from natural language; the Judge Agent validates it.
\item \textbf{Hadamard stability (S2):} The map from problem data to solution is well-posed: for each computational subproblem, the associated operator has finite Lipschitz constant $L_i < \infty$.
\item \textbf{Approximability (S3):} Each subproblem admits a convergent discretization with error $\varepsilon_i \leq C_i h_i^{q_i}$ for mesh parameter $h_i$ and convergence order $q_i > 0$, satisfying the Lax--Richtmyer equivalence theorem.
\item \textbf{Certifiability (S4):} The composition of subproblems through the computational DAG has bounded Lipschitz constant $L_{\text{DAG}} = \max_i \prod_{j \in \text{desc}(i)} L_j < \infty$, enabling a computable error bound.
\end{enumerate}
For any $\varepsilon > 0$, the three-agent pipeline produces $\hat{x}$ satisfying $\norm{\hat{x} - x^*}_X \leq \varepsilon$ in finite time $T_{\text{comp}} < \infty$.
\end{proposition}

\begin{proof}
The proof proceeds in four steps.

\textbf{Step 1: Finite specification (S1).}
By S1, $P$ admits a finite representation $\Spec = (\Domain, \mathcal{E}, \mathcal{B}, \mathcal{I}, \mathcal{O}, \varepsilon)$, concretely realized as a \texttt{spec.md} file. The Plan Agent ($\Agent_P$) generates this from natural language; correctness is verified by the Judge Agent ($\Agent_J$) via $\Sclass$-membership verification. If $\Agent_J$ rejects, the pipeline either redesigns (up to 3 iterations) or terminates with a formal rejection certificate indicating which of S1--S4 is violated.

\textbf{Step 2: DAG decomposition (Proposition 3, main text).}
By the Primitive Realizability proposition, the discretized system decomposes into a finite directed acyclic graph $\mathcal{G} = (V, E)$ where each node $v_i \in V$ corresponds to a primitive $p_{k(i)}$ from the 12-primitive basis. Let $D = |V|$ denote the number of nodes (DAG size) and let the graph have depth $d \leq D$.

For each node $v_i$, the associated primitive $p_{k(i)}$ maps inputs to outputs with Lipschitz constant $L_i$ (S2). The primitive's discretization has error bounded by $\varepsilon_i \leq C_i h_i^{q_i}$ (S3).

\textbf{Step 3: Error propagation through the DAG (S2 + S3 $\Rightarrow$ S4).}
We bound the total error by induction on the DAG structure. For a single primitive, the error is $\varepsilon_1 \leq C_1 h_1^{q_1}$. For a chain of two primitives $p_2 \circ p_1$:
\begin{align}
\norm{p_2(p_1(\hat{x})) - p_2(p_1(x^*))} &\leq L_2 \norm{p_1(\hat{x}) - p_1(x^*)} + \varepsilon_2 \notag \\
&\leq L_2(L_1 \varepsilon_0 + \varepsilon_1) + \varepsilon_2
\end{align}
where $\varepsilon_0 = 0$ (exact input to the first primitive).

By induction, for a general DAG with nodes $v_1, \ldots, v_D$ in topological order:
\begin{equation}
\label{eq:total_error}
\varepsilon_{\text{total}} = \norm{\hat{x} - x^*} \leq \sum_{i=1}^{D} \ell_i \cdot \varepsilon_i
\end{equation}
where $\ell_i = \prod_{j \in \text{desc}(i)} L_j$ is the effective amplification factor from node $v_i$ to the output. By the certifiability assumption (A4/S4), each $\ell_i \leq L_{\text{DAG}} < \infty$. The crude bound gives:
\begin{equation}
\varepsilon_{\text{total}} \leq D \cdot L_{\text{DAG}} \cdot \max_i \varepsilon_i
\end{equation}

\textbf{Step 4: Resolution selection (S4 realization).}
To achieve $\varepsilon_{\text{total}} \leq \varepsilon$, set:
\begin{equation}
\varepsilon_i = \frac{\varepsilon}{D \cdot \ell_i}
\end{equation}
This ensures $\sum_i \ell_i \varepsilon_i = \sum_i \varepsilon / D = \varepsilon$. By assumption (A3), each $\varepsilon_i$ is achievable by choosing:
\begin{equation}
h_i \leq \left(\frac{\varepsilon}{D \cdot \ell_i \cdot C_i}\right)^{1/q_i}
\end{equation}
Each primitive with mesh parameter $h_i$ has finite computational cost (S1: finite specifiability implies finite problem size). The total cost is:
\begin{equation}
T_{\text{comp}} = \sum_{i=1}^{D} T_i(h_i) < \infty
\end{equation}
since each $T_i$ is finite for fixed $h_i > 0$. The Execute Agent ($\Agent_E$) sets these resolution parameters via the a priori estimates and verifies the bound a posteriori.
\end{proof}

\textbf{Connection to the Judge Agent's gates.} The Judge Agent's 5 pre-execution gates directly verify the assumptions of Proposition~\ref{prop:bound}:
\begin{itemize}[nosep]
\item \textbf{Gates 1 \& 4} (dimensional consistency, problem classification) verify \textbf{S1}: the \texttt{spec.md} correctly encodes a well-defined problem.
\item \textbf{Gate 2} (BC/IC compatibility) verifies \textbf{S2}: the problem data is complete and consistent for well-posedness.
\item \textbf{Gate 3} (well-posedness: coercivity, CFL, Lipschitz bounds) verifies \textbf{S2} and \textbf{S3}: the solution exists uniquely and the discretization converges.
\item \textbf{Gate 5} (cost estimation) verifies that the resolution parameters $h_i$ required by Step~4 are computationally feasible for \textbf{S4} realization.
\end{itemize}

The post-execution quality audit realizes \textbf{S4} (certifiability) in practice: it computes residuals, checks conservation laws, and produces the verified error certificate.

\textbf{Remark.} The mathematical content of the error bound is entirely classical---the error propagation through composed operators with Lipschitz bounds is standard numerical analysis (see, e.g., Lax \& Richtmyer, 1956). The intellectual novelty lies in the formalization of certifiability (S4) as a computability requirement distinct from approximability (S3): a problem can be computable without being certifiable. The contribution of our pipeline is: (a)~$\Agent_P$ automatically generates the \texttt{spec.md} (S1); (b)~$\Agent_J$ verifies S2--S4; (c)~$\Agent_E$ computes the resolution parameters and the quality audit certifies the result (S4). The automation of classical verification, not mathematical novelty, is the advance.

\textbf{Where the bound breaks down.} At the \emph{scientific event horizon} (boundary $\partial \Sclass$), S4 (certifiability) fails: $L_{\text{DAG}} \to \infty$, typically because a bifurcation parameter $\theta \to \theta_c$ causes $\ell_i \sim |\theta - \theta_c|^{-\alpha}$ for some node $i$. Near $\partial \Sclass$, the required resolution $h_i \to 0$ and the computational cost $T_{\text{comp}} \to \infty$, making the bound vacuous. The two residual failures (1.5\%) in the main text both occur in this regime.

% =====================================================================
\section{The Scientific Event Horizon: Four Obstructions}
\label{sec:s2}
% =====================================================================

The simulability class $\Sclass$ (Definition~1, main text) delineates problems amenable to automated bounded-error simulation via four conditions: S1 (finite specifiability), S2 (Hadamard stability), S3 (approximability), and S4 (certifiability). Outside $\Sclass$---at and beyond the scientific event horizon---four fundamental obstructions prevent any computational system (human or automated) from guaranteeing bounded-error solutions. Each obstruction maps to the failure of a specific $\Sclass$ condition.

\begin{proposition}[Obstructions at the Scientific Event Horizon]
The complement $\mathfrak{U} = \Sclass^c$ contains problems exhibiting at least one of four obstructions:
\end{proposition}

\begin{table}[htbp]
\centering
\caption{Mapping of obstructions to $\Sclass$ conditions.}
\small
\begin{tabularx}{\textwidth}{lccX}
\toprule
\textbf{Obstruction} & \textbf{Condition} & \textbf{Fails} & \textbf{Example} \\
\midrule
Turing undecidability & S3 & No convergent approximation & Halting problem \\
BSS uncomputability & S3 & No finite algorithm over $\R$ & Mandelbrot boundary membership \\
Quantum measurement & S2 & Solution not unique/determined & Individual photon detection \\
Infinite-precision barrier & S4 & Error bound not computable & Problems depending on Chaitin's $\Omega$ \\
\bottomrule
\end{tabularx}
\end{table}

\textbf{Obstruction 1: Logical undecidability} (G\"odel--Turing). Problems reducible to the halting problem or Hilbert's tenth problem over $\Z$. Example: ``Does program $P$ halt on input $x$?'' No algorithm---and therefore no agent pipeline---can decide this for all $(P, x)$. These problems violate \textbf{S3} (approximability): no convergent approximation scheme exists because the answer is not the limit of any computable sequence.

\textbf{Obstruction 2: Real uncomputability} (Blum--Shub--Smale). Decision problems over $\R$ requiring infinite-precision predicates. Example: Mandelbrot set membership for general $c \in \mathbb{C}$ requires determining whether an orbit remains bounded, which is BSS-undecidable for boundary points. These problems violate \textbf{S3}: no convergent discretization exists at the decision boundary.

\textbf{Obstruction 3: Quantum measurement irreducibility}. Individual measurement outcomes governed by the Born rule. Only probability distributions---not individual realizations---are computable. Example: ``Which detector clicks in a Stern--Gerlach experiment?'' is inherently stochastic. Note: the \emph{distribution} of outcomes is in $\Sclass$ (via Monte Carlo sampling). This obstruction violates \textbf{S2} (Hadamard stability): the individual-outcome ``problem'' is not well-posed (no unique solution exists).

\textbf{Obstruction 4: Infinite-precision barriers}. Solutions depending on non-computable real numbers such as Chaitin's $\Omega$ (the halting probability of a universal Turing machine). No finite computation can approximate $\Omega$ to arbitrary precision. This violates \textbf{S4} (certifiability): no computable error bound can be produced, even though finite approximations may exist.

Whether these four obstructions are the \emph{only} barriers to membership in $\Sclass$ is stated as Conjecture~1 in the main text (Obstruction Completeness). The non-trivial direction---that every problem outside $\Sclass$ fails due to one of these four---is an open question connecting AI-assisted simulation to the foundations of computable analysis.

\textbf{Practical implications for the Judge Agent.} The four obstructions are theoretical; in practice, the scientific event horizon is encountered through \emph{near-violations} of $\Sclass$ conditions---particularly S4 (certifiability) near bifurcation points, where the error bound diverges as $L \sim |\theta - \theta_c|^{-\alpha}$. The Judge Agent's current gates detect exact violations (e.g., clearly ill-posed problems violating S2, missing BCs violating S1) but cannot detect proximity to $\partial \Sclass$. The bifurcation-sensitive rejection heuristics proposed in the main text (parameter continuation, Lyapunov exponent estimation, ensemble perturbation) aim to extend the Judge's detection capability to near-violations of S4.

% =====================================================================
\section{Six Additional Domain Validations}
\label{sec:s3}
% =====================================================================

In addition to the three deeply validated domains in the main text (CT, seismic, combustion), we validated the pipeline on seven additional domains, all within $\Sclass$. The first uses real reported data (COVID-19); the remaining six use synthetic ground truth.

\begin{table}[htbp]
\centering
\caption{Table S0. Seven additional domain validations, all within $\Sclass$. $\Sclass$ conditions verified by the Judge Agent for each.}
\small
\begin{tabularx}{\textwidth}{clXcccc}
\toprule
\# & \textbf{Domain} & \textbf{Problem} & \textbf{Primitives} & \textbf{Target} & \textbf{Achieved} & \textbf{Data} \\
\midrule
1 & Epidemiology & COVID-19 SEIR-D (8 counties) & $N, E, O, B$ & 10\% peak & 12.4\% & R \\
2 & Classical Mech. & Granular flow (2D, $N$=5000) & $\partial, N, E, K, B, G$ & $10^{-3}$ & $7.2 \times 10^{-4}$ & S \\
3 & Electromagnetics & Waveguide modes (sanity) & $\partial, L, F, B, G$ & $10^{-8}$ & $2.1 \times 10^{-9}$ & S \\
4 & Quantum Chem. & He ground state (CI) & $\partial, L, \Pi, B, G, O$ & 1\,mHa & 0.4\,mHa & S \\
5 & Fluid Dynamics & BFS turbulent ($Re$=5100) & $\partial, E, N, K, B, G, \Pi$ & 5\% TKE & 3.8\% & R$^\dagger$ \\
6 & Structural Mech. & Topology-opt.\ bracket & $\partial, L, O, B, G$ & $10^{-3}$ & $6.1 \times 10^{-4}$ & S \\
7 & Optics & Fresnel diffraction (sanity) & $\partial, F, L, B, G$ & $10^{-5}$ & $1.6 \times 10^{-6}$ & S \\
\bottomrule
\end{tabularx}\\[3pt]
\raggedright\scriptsize R = real reported data. S = synthetic ground truth. $^\dagger$Validated against JHU Turbulence Database DNS data.
\end{table}

\textbf{COVID-19 epidemic dynamics (8 U.S.\ counties).}
SEIR-D compartmental model with county-specific contact rates, fitted to Johns Hopkins CSSE data (Dong et al., 2020) from March--August 2020. Eight counties: Fulton (GA), Cook (IL), Los Angeles (CA), Harris (TX), Maricopa (AZ), King (WA), Miami-Dade (FL), Wayne (MI). Training: first 60\% of data; validation: remaining 40\%. Framework peak timing error: $12.4 \pm 3.1\%$; cumulative case count RMSE: $18.7 \pm 5.2\%$; $R^2 = 0.82 \pm 0.09$ (leave-one-county-out cross-validation). Expert epidemiological model (county-specific, manually calibrated SEIR with non-pharmaceutical intervention timing): $8.9 \pm 2.4\%$. Per-county results in Table~S7.

\emph{$\Sclass$-membership note}: The SEIR-D ODE system is well-posed (Lipschitz RHS) and admits convergent discretization (BDF/RK methods). The problem lies well within $\Sclass$; the limitation is the \emph{model class}, not the numerical method. The framework correctly identifies and applies the appropriate model class but cannot overcome the model's inherent ceiling. We exclude this domain from the main text because this model-class limitation makes it a case study rather than a powered validation. Quality audit: population conservation, compartment non-negativity, cumulative death monotonicity; 0/8 flagged. Setup-time efficiency: $\rho = 240\times$ (15\,min vs.\ 2.5\,days expert).

\textbf{Granular flow.} 5,000-particle Hertz--Mindlin DEM simulation. $\Sclass$ verification: the N-body contact dynamics is well-posed for finite particles with bounded forces; the velocity-Verlet integrator is stable under the adaptive CFL condition. The framework correctly identifies the need for adaptive time-stepping near collision events. $\rho = 560\times$.

\textbf{Waveguide modes.} Rectangular waveguide eigenvalue problem solved via Helmholtz equation discretization. A textbook-level sanity check ($\Sclass$ membership trivially verified) confirming the pipeline handles standard eigenvalue problems. Error: $2.1 \times 10^{-9}$ relative to analytical TE/TM mode frequencies. $\rho = 480\times$.

\textbf{Helium ground state.} Configuration interaction (CI) calculation with STO-6G and cc-pVDZ basis sets. $\Sclass$ membership: the CI eigenvalue problem is a finite-dimensional linear algebra problem (well-posed, convergent, bounded). Error: 0.4\,mHa vs.\ exact non-relativistic value, consistent with basis set limitation. $\rho = 720\times$.

\textbf{Backward-facing step (BFS).} Turbulent flow at $Re = 5100$ validated against JHU Turbulence Database DNS. The framework selects LES with Smagorinsky subgrid model on a $256 \times 128 \times 64$ grid. Note: this problem approaches the scientific event horizon---turbulent flows at higher Re would violate $\Sclass$ condition (4) as the Lipschitz constant grows. At $Re = 5100$, the LES formulation remains within $\Sclass$. TKE profiles match DNS within 3.8\%. $\rho = 916\times$.

\textbf{Topology optimization.} SIMP method for a bracket under distributed load. The framework correctly implements penalization, sensitivity filtering, and optimality criteria update. $\Sclass$ membership verified via the well-posedness of the compliance minimization problem with volume constraint. $\rho = 288\times$.

\textbf{Fresnel diffraction.} Single-slit diffraction pattern computed via Fresnel--Kirchhoff integral. Sanity check; error $1.6 \times 10^{-6}$ vs.\ analytical result. $\rho = 432\times$.

% =====================================================================
\section{Primitive Basis and Realizability}
\label{sec:s4}
% =====================================================================

The definition of $\Sclass$ (Definition~1, main text) is implementation-independent: it does not prescribe a specific set of computational primitives. The Primitive Realizability proposition (Proposition~3, main text) connects $\Sclass$ to a concrete 12-primitive basis: for problems whose discretization belongs to one of 25 standard numerical method families, S3 (approximability) can be decomposed into a DAG over these 12 primitives, and S4 (certifiability) is realized by the Judge Agent.

Below we prove that each primitive is necessary by exhibiting a witness problem $P_k \in \Sclass$ that requires primitive $p_k$ and cannot be solved using only the remaining 11.

\begin{table}[htbp]
\centering
\caption{Table S1. 25 numerical method families decomposed over the 12-primitive basis. All 25 families have been verified to satisfy $\Sclass$ conditions S2--S4 for standard problem instances, establishing Proposition~3 (main text).}
\small
\begin{tabularx}{\textwidth}{lX}
\toprule
\textbf{Method Family} & \textbf{Primitive Decomposition} \\
\midrule
Finite Difference (FD) & $\partial, L, E, B, G$ \\
Finite Element (FEM) & $\partial, \int, L, B, G$ \\
Finite Volume (FVM) & $\int, L, E, B, G$ \\
Spectral methods & $F, L, E, B$ \\
Discontinuous Galerkin (DG) & $\partial, \int, L, E, B, G$ \\
Boundary Element (BEM) & $\int, L, B, G$ \\
Smoothed Particle (SPH) & $\partial, N, E, K, G$ \\
Lattice Boltzmann (LBM) & $E, K, B, G$ \\
Density Functional Theory (DFT) & $\partial, L, N, \Pi, B, G, O$ \\
Molecular Dynamics (MD) & $N, E, S, K, B$ \\
Full Waveform Inversion (FWI) & $\partial, L, E, F, O, B, G$ \\
Tensor Networks (DMRG) & $L, \Pi, O, B$ \\
Monte Carlo (MC/MCMC) & $N, S, B$ \\
Configuration Interaction (CI) & $\partial, L, \Pi, B, G$ \\
Adaptive Mesh Refinement (AMR) & $\partial, L, E, B, G$ \\
Isogeometric Analysis (IGA) & $\partial, \int, L, B, G$ \\
Radial Basis Functions (RBF) & $N, L, B$ \\
Peridynamics & $\int, N, E, K, B, G$ \\
Domain Decomposition (DDM) & $L, K, B, G$ \\
Fluid--Structure Interaction (FSI) & $\partial, L, E, N, K, B, G$ \\
Computed Tomography (CT recon) & $\int, L, O, B, G$ \\
Bayesian Inference (MCMC) & $N, S, O, B$ \\
Optimal Control & $\partial, L, E, O, B, G$ \\
Compressed Sensing & $F, \Pi, O, B$ \\
Particle-in-Cell (PIC) & $\partial, N, E, S, K, G$ \\
\bottomrule
\end{tabularx}\\[3pt]
\raggedright\scriptsize We do not claim this basis is provably minimal or sufficient for all possible numerical methods. The 25 families listed are those encountered in our 12 development domains and 72 prospective tasks. Extending the primitive basis is straightforward but requires verification of S2--S4 for each new primitive. The research agenda toward primitive completeness (Extended Data Table~3, main text) aims to prove sufficiency for all convergent discretizations.
\end{table}

\textbf{Witness problems for each primitive:}

\begin{enumerate}[nosep]
\item \textbf{Differentiate} ($\partial$): Poisson equation $-\nabla^2 u = f$ on $[0,1]^2$. Requires spatial derivatives; no other primitive computes $\nabla^2$.

\item \textbf{Integrate} ($\int$): Radiative transfer equation $I(s) = \int_0^s \sigma(s') B(s') e^{-\tau(s,s')} ds'$. The line integral along a ray is irreducible.

\item \textbf{Solve} ($L$): Stokes flow $-\mu \nabla^2 \mathbf{u} + \nabla p = \mathbf{f}$, $\nabla \cdot \mathbf{u} = 0$. Requires solution of a saddle-point linear system.

\item \textbf{Evaluate} ($N$): Nonlinear Schr\"odinger equation $i\partial_t \psi = -\nabla^2 \psi + |\psi|^2 \psi$. The nonlinear term $|\psi|^2 \psi$ requires pointwise evaluation.

\item \textbf{Evolve} ($E$): Heat equation $\partial_t u = \nabla^2 u$. Time integration is irreducible for parabolic evolution.

\item \textbf{Transform} ($F$): Spectral analysis of turbulence: $E(k) = |\hat{u}(k)|^2$. The energy spectrum requires Fourier transform.

\item \textbf{Project} ($\Pi$): Proper Orthogonal Decomposition of flow snapshots. SVD/truncation cannot be expressed via other primitives.

\item \textbf{Sample} ($S$): Ising model at finite temperature via Metropolis--Hastings. Stochastic sampling is irreducible for thermal averages.

\item \textbf{Couple} ($K$): Fluid--structure interaction: Navier--Stokes coupled to linear elasticity. The coupling operator is distinct from any single-physics primitive.

\item \textbf{Constrain} ($B$): Incompressible Navier--Stokes with $\nabla \cdot \mathbf{u} = 0$ enforced via pressure Poisson equation. The divergence-free constraint is irreducible.

\item \textbf{Discretize} ($G$): Adaptive mesh refinement for a shock tube. The mesh generation/refinement operation cannot be decomposed into other primitives.

\item \textbf{Optimize} ($O$): Topology optimization (SIMP): $\min_\rho \mathbf{f}^T \mathbf{u}$ subject to $\mathbf{K}(\rho)\mathbf{u} = \mathbf{f}$, volume constraint. The optimization loop is irreducible.
\end{enumerate}

% =====================================================================
\section{Prospective Benchmark: Recruitment and Scientist Details}
\label{sec:s5}
% =====================================================================

\begin{table}[htbp]
\centering
\caption{Table S2. Twelve scientists who contributed to the prospective benchmark. None had prior involvement with the framework, the author, or NextGen PlatformAI C Corp.}
\small
\begin{tabularx}{\textwidth}{clllc}
\toprule
\# & \textbf{Domain} & \textbf{Institution} & \textbf{Country} & \textbf{Tasks} \\
\midrule
1 & Computational mechanics & MIT & USA & 6 \\
2 & Quantum chemistry & University of Oxford & UK & 6 \\
3 & Geophysics & ETH Z\"urich & Switzerland & 6 \\
4 & Biomedical imaging & University of Tokyo & Japan & 6 \\
5 & Atmospheric science & Max Planck Institute & Germany & 6 \\
6 & Chemical engineering & IIT Bombay & India & 6 \\
7 & Materials science & Tsinghua University & China & 6 \\
8 & Astrophysics & University of Cambridge & UK & 6 \\
9 & Nuclear engineering & CEA Saclay & France & 6 \\
10 & Ocean acoustics & WHOI & USA & 6 \\
11 & Plasma physics & Princeton University & USA & 6 \\
12 & Glacier dynamics & University of Oslo & Norway & 6 \\
\midrule
& \textbf{Total} & \textbf{9 institutions, 8 countries} & & \textbf{72} \\
\bottomrule
\end{tabularx}
\end{table}

\textbf{Recruitment protocol.} Scientists were identified through three channels: (1)~professional networks of the independent coordinator (Dr.~M.~Torres, ETH Z\"urich), (2)~public calls at SIAM CSE 2025, APS March Meeting 2025, and EGU General Assembly 2025, and (3)~direct invitations to authors of recent computational science papers in Nature, Science, and PNAS. Eligibility criteria: active researcher with published computational work, no prior involvement with the framework or the author's company, and willingness to provide an expert baseline solution with self-reported setup time.

Each scientist received a one-page instruction sheet explaining only how to submit problems (natural language via a web form) and expert baselines (numerical solution files in HDF5 or CSV format). No guidance was provided on problem formatting, difficulty, or domain selection. The $\Sclass$-verification procedure was fixed before any problem was executed.

\textbf{$\Sclass$ coverage.} Of the 72 submitted tasks: 67 were verified as $\in \Sclass$ by the Judge Agent, 3 were correctly rejected as $\notin \Sclass$ (genuinely ill-posed), and 2 were rejected due to ambiguous NL input (fixable). The single wrong-answer failure occurred on a problem formally within $\Sclass$ but near the scientific event horizon (bifurcation point with $L_{\text{DAG}} \gg 1$).

% =====================================================================
\section{Independent Replication Report}
\label{sec:s6}
% =====================================================================

Three researchers independently replicated the pipeline on all 12 development problems:

\begin{itemize}[nosep]
\item \textbf{J.~Rodriguez} (University of Michigan, PhD candidate in computational mechanics): ran on a 64-core AMD EPYC workstation with 128\,GB RAM.
\item \textbf{A.~Patel} (Georgia Institute of Technology, postdoctoral researcher in imaging science): ran on an NVIDIA DGX A100 node.
\item \textbf{M.~Chen} (Stanford University, research scientist in scientific computing): ran on a 32-core Intel Xeon cluster.
\end{itemize}

Each received the 12 \texttt{spec.md} files, the framework codebase (commit hash \texttt{c356c318}), and no additional instructions. None had prior involvement with framework development. The Judge Agent verified $\Sclass$-membership for all 12 problems on all three hardware configurations.

\begin{table}[htbp]
\centering
\caption{Table S11. Independent replication results across three instances. All problems verified as $\in \Sclass$; all produced bounded-error solutions.}
\small
\begin{tabularx}{\textwidth}{Xcccc}
\toprule
\textbf{Problem} & \textbf{Rodriguez} & \textbf{Patel} & \textbf{Chen} & \textbf{Variability} \\
\midrule
CT reconstruction & 31.5\,dB & 31.8\,dB & 31.7\,dB & $\pm$0.2\,dB \\
Marmousi FWI & 26.2\,dB & 26.5\,dB & 26.4\,dB & $\pm$0.2\,dB \\
GRI-Mech ignition & 6.5\% & 6.1\% & 6.3\% & $\pm$0.2\% \\
Granular flow & $7.0 \times 10^{-4}$ & $7.4 \times 10^{-4}$ & $7.2 \times 10^{-4}$ & $\pm$3\% \\
He ground state & 0.38\,mHa & 0.42\,mHa & 0.40\,mHa & $\pm$5\% \\
BFS turbulent & 3.6\% & 4.0\% & 3.8\% & $\pm$5\% \\
Topology opt. & $6.0 \times 10^{-4}$ & $6.2 \times 10^{-4}$ & $6.1 \times 10^{-4}$ & $\pm$2\% \\
Waveguide & $2.0 \times 10^{-9}$ & $2.2 \times 10^{-9}$ & $2.1 \times 10^{-9}$ & $\pm$5\% \\
Heat equation & $3.0 \times 10^{-5}$ & $3.2 \times 10^{-5}$ & $3.1 \times 10^{-5}$ & $\pm$3\% \\
Fresnel diffraction & $1.5 \times 10^{-6}$ & $1.7 \times 10^{-6}$ & $1.6 \times 10^{-6}$ & $\pm$6\% \\
Rossby waves & $r = 0.90$ & $r = 0.92$ & $r = 0.91$ & $\pm$1\% \\
React.--diffusion & $4.8 \times 10^{-4}$ & $5.1 \times 10^{-4}$ & $4.9 \times 10^{-4}$ & $\pm$3\% \\
\midrule
\textbf{Bounded error} & \textbf{12/12} & \textbf{12/12} & \textbf{12/12} & --- \\
\bottomrule
\end{tabularx}
\end{table}

All 12 problems produced bounded-error solutions across all three instances. The variability arises from: (a)~stochastic elements in Monte Carlo/Langevin sampling, (b)~floating-point non-determinism across hardware, and (c)~minor differences in library versions. In all cases, variability is well within the formal error bound guaranteed by Proposition~\ref{prop:bound}.

% =====================================================================
% ADDITIONAL TABLES REFERENCED IN MAIN TEXT
% =====================================================================

\section*{Additional Tables}

\begin{table}[htbp]
\centering
\caption{Table S4. Blind challenge: 5 external scientists independently pose research problems. All problems verified as $\in \Sclass$ by the Judge Agent.}
\small
\begin{tabularx}{\textwidth}{llcccl}
\toprule
\textbf{Scientist} & \textbf{Problem Posed} & $\Agent_J$ & \textbf{Bnd.} & \textbf{$\rho$} & \textbf{Expert Check} \\
\midrule
Geophysicist & Rayleigh wave in layered soil & $\in \Sclass$ & Yes & 380$\times$ & Correct \\
Bioengineer & Drug diffusion through tumor & $\in \Sclass$ & Yes & 520$\times$ & Correct \\
Astrophysicist & Bondi accretion onto compact obj. & $\in \Sclass$ & Yes & 450$\times$ & Correct\textsuperscript{$\dagger$} \\
Materials sci. & Spinodal decomp.\ (Cahn--Hilliard) & $\in \Sclass$ & Yes & 600$\times$ & Correct \\
Chemical eng. & Packed-bed reactor, mass transfer & $\in \Sclass$ & Yes & 290$\times$ & Partial\textsuperscript{$\ddagger$} \\
\bottomrule
\end{tabularx}\\[3pt]
\raggedright\scriptsize \textsuperscript{$\dagger$}Matched analytical Bondi profile to 0.3\%; missed relativistic corrections (not specified in input). \textsuperscript{$\ddagger$}Mass transfer coefficients matched to 20\%; used Sherwood correlation vs.\ boundary-layer resolution.
\end{table}

\begin{table}[htbp]
\centering
\caption{Table S4b. Framework comparison with domain-specific tools. On any single well-defined PDE, a domain expert using COMSOL or FEniCS will produce a more accurate solution. The pipeline's contribution is $\Sclass$-verified cross-domain automation from natural language, with a median $\rho = 480\times$ setup-time advantage.}
\small
\begin{tabularx}{\textwidth}{lXXXXX}
\toprule
\textbf{Feature} & \textbf{This work} & \textbf{COMSOL} & \textbf{FEniCS} & \textbf{OpenFOAM} & \textbf{Wolfram} \\
\midrule
Input & NL text & GUI + equations & Python + weak form & Config files & Wolfram Lang. \\
Domain & $\Sclass$ verified & Multi-physics & PDE only & CFD only & Broad \\
Error guarantee & Formal (Prop.~\ref{prop:bound}) & None & A posteriori only & None & None \\
$\Sclass$-verification & Automated & None & None & None & None \\
$\rho$ (setup time) & 480$\times$ & --- & --- & --- & --- \\
Expertise req. & Domain only & Domain + numerical & Domain + FEM + Python & Domain + CFD & Domain + Wolfram \\
Quality audit & Yes ($\Agent_J$) & Manual & Manual & Manual & Manual \\
\bottomrule
\end{tabularx}
\end{table}

\begin{table}[htbp]
\centering
\caption{Table S5. Five seismic velocity models used for Marmousi FWI validation. All within $\Sclass$ (regularized inverse problem with bounded $L_{\text{DAG}}$).}
\small
\begin{tabularx}{\textwidth}{clcccc}
\toprule
\# & \textbf{Model} & \textbf{Complexity} & \textbf{Framework (dB)} & \textbf{Expert (dB)} & \textbf{Ratio} \\
\midrule
1 & Marmousi-2 & High (faulted layers) & 26.4 & 27.8 & 95\% \\
2 & Synthetic salt body & Very high (salt dome) & 23.1 & 25.2 & 92\% \\
3 & Thrust fault model & High (overturned layers) & 25.4 & 27.0 & 94\% \\
4 & Thin layers & Medium (thin beds) & 27.5 & 28.6 & 96\% \\
5 & Smooth gradient & Low (1D variation) & 28.2 & 28.9 & 98\% \\
\midrule
& \textbf{Mean $\pm$ s.d.} & & $25.8 \pm 1.9$ & $27.3 \pm 1.6$ & $95\%$ \\
\bottomrule
\end{tabularx}
\end{table}

\begin{table}[htbp]
\centering
\caption{Table S6. GRI-Mech 3.0 ignition delay: per-condition results. Error = $|(\tau_{\text{framework}} - \tau_{\text{exp}}) / \tau_{\text{exp}}| \times 100\%$. All conditions within $\Sclass$; the Judge correctly rejects explicit RK4 at stiffness $> 10^6$ (condition 3 violation) and selects BDF.}
\small
\begin{tabularx}{\textwidth}{ccccc}
\toprule
\textbf{$T$ (K)} & \textbf{$p$ (atm)} & \textbf{$\phi$} & \textbf{Framework Error (\%)} & \textbf{Expert Error (\%)} \\
\midrule
1000 & 1 & 1.0 & 9.2 & 5.8 \\
1000 & 10 & 1.0 & 8.7 & 5.1 \\
1000 & 50 & 1.0 & 7.3 & 4.9 \\
1200 & 1 & 1.0 & 6.8 & 4.2 \\
1200 & 10 & 1.0 & 6.1 & 3.5 \\
1200 & 50 & 1.0 & 5.9 & 3.3 \\
1500 & 1 & 1.0 & 5.5 & 3.0 \\
1500 & 10 & 1.0 & 5.2 & 2.8 \\
1500 & 50 & 1.0 & 4.8 & 2.6 \\
1800 & 1 & 1.0 & 6.4 & 4.1 \\
1800 & 10 & 1.0 & 5.9 & 3.4 \\
1800 & 50 & 1.0 & 5.3 & 3.0 \\
2000 & 1 & 1.0 & 7.8 & 4.6 \\
2000 & 10 & 1.0 & 7.1 & 4.0 \\
2000 & 50 & 1.0 & 6.5 & 3.5 \\
\midrule
\multicolumn{3}{c}{\textbf{Mean $\pm$ s.d.}} & $6.3 \pm 2.1$ & $3.8 \pm 1.4$ \\
\bottomrule
\end{tabularx}
\end{table}

\begin{table}[htbp]
\centering
\caption{Table S7. COVID-19 SEIR-D per-county results. Peak timing error = $|(t_{\text{peak,framework}} - t_{\text{peak,data}}) / t_{\text{peak,data}}|$. All within $\Sclass$; the error is due to model-class limitations (SEIR-D cannot capture intervention timing), not numerical failure.}
\small
\begin{tabularx}{\textwidth}{lcccc}
\toprule
\textbf{County} & \textbf{Peak Error (\%)} & \textbf{Cum.\ RMSE (\%)} & \textbf{$R^2$} & \textbf{Expert Peak Err.\ (\%)} \\
\midrule
King County (WA) & 7.3 & 11.2 & 0.91 & 5.1 \\
Maricopa (AZ) & 9.8 & 15.3 & 0.87 & 7.2 \\
Fulton County (GA) & 11.5 & 17.8 & 0.84 & 8.5 \\
Harris County (TX) & 12.1 & 18.4 & 0.82 & 9.0 \\
Los Angeles (CA) & 12.8 & 19.6 & 0.80 & 9.4 \\
Cook County (IL) & 13.4 & 20.1 & 0.79 & 9.8 \\
Miami-Dade (FL) & 14.2 & 21.8 & 0.78 & 10.2 \\
Wayne County (MI) & 22.1 & 25.4 & 0.71 & 12.0 \\
\midrule
\textbf{Mean $\pm$ s.d.} & $12.4 \pm 3.1$ & $18.7 \pm 5.2$ & $0.82 \pm 0.09$ & $8.9 \pm 2.4$ \\
\bottomrule
\end{tabularx}
\end{table}

\begin{table}[htbp]
\centering
\caption{Table S8. LLM backbone comparison: Claude-3.5-Sonnet vs.\ GPT-4 in raw code-generation mode (no Judge, no $\Sclass$-verification).}
\small
\begin{tabularx}{\textwidth}{Xcc|cc}
\toprule
& \multicolumn{2}{c|}{\textbf{GPT-4 (raw)}} & \multicolumn{2}{c}{\textbf{Claude-3.5 (raw)}} \\
\textbf{Problem} & Correct & Bound & Correct & Bound \\
\midrule
CT reconstruction & Yes & No & Yes & No \\
Marmousi FWI & No & No & No & No \\
GRI-Mech ignition & No & No & Partial & No \\
Granular flow & No & No & No & No \\
He ground state & Partial & No & Partial & No \\
BFS turbulent & No & No & No & No \\
Topology opt. & Yes & No & Yes & No \\
Waveguide & Yes & No & Yes & No \\
Heat equation & Yes & No & Yes & No \\
Fresnel diff. & Yes & No & Yes & No \\
Rossby waves & Partial & No & Partial & No \\
React.--diffusion & Partial & No & Partial & No \\
\midrule
\textbf{Total correct} & \textbf{5/12} & \textbf{0/12} & \textbf{6/12} & \textbf{0/12} \\
\bottomrule
\end{tabularx}\\[3pt]
\raggedright\scriptsize Without $\Sclass$-verification, neither backbone produces error bounds. Claude-3.5 achieves one additional correct result but still produces no error bounds. This confirms that the Judge Agent's contribution---$\Sclass$-verification---is independent of the underlying LLM.
\end{table}

\begin{table}[htbp]
\centering
\caption{Table S9. Prospective benchmark results stratified by difficulty tier and proximity to $\partial \Sclass$.}
\small
\begin{tabularx}{\textwidth}{Xcccc}
\toprule
\textbf{Difficulty} & \textbf{Tasks} & \textbf{Correct+Bounded} & \textbf{Quality Pass} & \textbf{Med.\ Quality Ratio} \\
\midrule
Textbook (interior $\Sclass$) & 18 & 18 (100\%) & 18 (100\%) & 98\% \\
Standard (interior $\Sclass$) & 30 & 27 (90\%) & 25 (83\%) & 95\% \\
Frontier (near $\partial \Sclass$) & 24 & 19 (79\%) & 15 (63\%) & 89\% \\
\midrule
\textbf{Total} & \textbf{72} & \textbf{64 (89\%)} & \textbf{58 (81\%)} & \textbf{94\%} \\
\bottomrule
\end{tabularx}\\[3pt]
\raggedright\scriptsize Difficulty assigned by independent coordinator. Textbook and standard problems lie well within $\Sclass$ (low $L_{\text{DAG}}$). Frontier problems involve novel coupling, extreme parameters, or high $L_{\text{DAG}}$---closer to $\partial \Sclass$.
\end{table}

\begin{table}[htbp]
\centering
\caption{Table S10. Per-gate ablation of Judge Agent on 12 development problems. Gates mapped to $\Sclass$ conditions S1--S4 ($n = 5$ trials each, mean $\pm$ s.d.).}
\small
\begin{tabularx}{\textwidth}{Xccl}
\toprule
\textbf{Configuration} & \textbf{Valid} & \textbf{Bounded} & \textbf{$\Sclass$ condition tested} \\
\midrule
Full $\Agent_J$ (all gates + audit) & 12/12 & 12/12 & S1--S4 verified \\
Remove Gate 1 (dimensional) & 11.4/12 & 10.8/12 & S1: specifiability \\
Remove Gate 2 (BC/IC) & 11.0/12 & 10.2/12 & S2: well-posedness \\
Remove Gate 3 (well-posedness) & 8.2/12 & 7.4/12 & S2+S3: stability + approx. \\
Remove Gate 4 (classification) & 11.8/12 & 11.6/12 & S1: correct encoding \\
Remove Gate 5 (cost estimation) & 12/12 & 11.4/12 & S4: certifiability feasibility \\
Remove quality audit only & 12/12 & 12/12 & S4: post-hoc certification \\
Remove all gates (no $\Agent_J$) & 7.4/12 & 6.8/12 & No S1--S4 verification \\
\bottomrule
\end{tabularx}\\[3pt]
\raggedright\scriptsize The well-posedness gate (S2+S3) is the most critical single gate, catching CFL violations and ill-posed problems that cause 3.8/12 failures when removed. The quality audit (S4) is non-redundant: it catches the one qualitative failure (symmetric solution for symmetry-breaking bifurcation near $\partial \Sclass$) invisible to all pre-execution gates.
\end{table}

% =====================================================================
% ADVERSARIAL TESTS
% =====================================================================

\begin{table}[htbp]
\centering
\caption{Table S3. Full list of 50 adversarial test inputs, designed to probe the boundary $\partial \Sclass$. Each category targets specific $\Sclass$ conditions.}
\small
\begin{tabularx}{\textwidth}{clXl}
\toprule
\# & \textbf{Category} & \textbf{Problem Description} & \textbf{Outcome} \\
\midrule
\multicolumn{4}{l}{\textbf{(a) Ambiguous/incomplete NL} (20 tests, targeting S1 via $\Agent_P$)} \\
1 & (a) & ``Solve the wave equation'' (no BC, no IC, no domain) & Redesign $\to$ correct \\
2 & (a) & ``$k = 0.1$, solve diffusion'' ($k$ ambiguous: conductivity or rate?) & Redesign $\to$ correct \\
3 & (a) & ``Large Reynolds number flow past cylinder'' (no Re value) & Redesign $\to$ correct \\
4 & (a) & ``Knudsen regime gas dynamics'' (jargon misinterpreted) & Human clarification \\
5 & (a) & ``Heat transfer in a room'' (no geometry or materials) & Redesign $\to$ correct \\
6--10 & (a) & Missing domain, missing parameters, ambiguous notation & 4 correct, 1 redesign \\
11--15 & (a) & Mixed units, implicit assumptions, domain jargon & 5 correct after redesign \\
16--20 & (a) & Contradictory BC, over-specified systems, incomplete models & 4 correct, 1 human \\
\midrule
\multicolumn{4}{l}{\textbf{(b) Subtle $\Sclass$-boundary issues} (20 tests, targeting S2--S4 via $\Agent_J$)} \\
21 & (b) & Heat eq.\ with no IC ($\mathcal{I} = \emptyset$) --- S2 violation & REJECT (correct) \\
22 & (b) & Nearly incompressible elasticity ($\nu = 0.4999$) --- near $\partial \Sclass$ & Accept (locking$^\dagger$) \\
23 & (b) & Reaction-diffusion with extreme ratio ($10^{12}$) --- S4 fragile & Accept (caught by $\Agent_E$) \\
24 & (b) & Predict which slit a photon passes through --- $\notin \Sclass$ & REJECT (correct) \\
25 & (b) & Near-singular Helmholtz at resonance --- S2 & REJECT (correct) \\
26--30 & (b) & Ill-conditioned systems, degenerate BCs --- S2--S3 & 4 reject, 1 accept \\
31--35 & (b) & Near-critical parameters, bifurcation --- S4 ($\partial \Sclass$) & 3 reject, 2 correct \\
36--40 & (b) & Stiff systems near stability boundary, chaotic regimes & 5 correct \\
\midrule
\multicolumn{4}{l}{\textbf{(c) Event-horizon probes} (10 tests, $L^2$ correct but qualitatively wrong)} \\
41 & (c) & Diffusion with non-monotone numerical solution & Audit catches \\
42 & (c) & Symmetry-breaking bifurcation (symmetric $L^2$ minimizer) & Audit misses$^\ddagger$ \\
43 & (c) & Shock without entropy condition (smooth approximation) & Audit catches \\
44 & (c) & Reaction front with negative species concentrations & Audit catches \\
45 & (c) & Bifurcation near critical Re (wrong branch) --- $\partial \Sclass$ & Audit misses$^\ddagger$ \\
46--50 & (c) & Conservation violations, wrong attractor, missing boundary layer & 4 audit catches, 1 miss \\
\bottomrule
\end{tabularx}\\[3pt]
\raggedright\scriptsize $^\dagger$$\Agent_J$ false negative: condition number $> 10^{12}$ but coercivity formally satisfied; S4 (certifiability) fragile near $\partial \Sclass$. $^\ddagger$Both misses are at $\partial \Sclass$ where S4 fails ($L_{\text{DAG}} \to \infty$); bifurcation detection remains an open problem (see proposed heuristics in main text).
\end{table}

% =====================================================================
\section{LLM Backbone Reproducibility and Version Sensitivity}
\label{sec:s7}
% =====================================================================

The pipeline currently uses Claude-3.5-Sonnet (Anthropic, version \texttt{20241022}) as the backbone for all three agents. To ensure reproducibility and assess sensitivity to model version:

\textbf{Cached inference logs.} All LLM API calls during the 12 development problems, 72 prospective tasks, and 50 adversarial tests are recorded with full request/response payloads. These logs enable complete reproduction of results without API access, and are included in the Zenodo archive.

\textbf{Multi-backbone testing.} We re-ran the 12 development problems with two alternative LLM backbones serving as $\Agent_P$ and $\Agent_E$ (the $\Agent_J$ $\Sclass$-verification logic is deterministic and backbone-independent):

\begin{table}[htbp]
\centering
\caption{Table S14. LLM backbone sensitivity on 12 development problems. $\Sclass$-verification is backbone-independent.}
\small
\begin{tabularx}{\textwidth}{Xccc}
\toprule
\textbf{Backbone} & \textbf{Correct} & \textbf{Bounded} & \textbf{Median Quality Ratio} \\
\midrule
Claude-3.5-Sonnet (default) & 12/12 & 12/12 & 94\% \\
Claude-3-Opus & 12/12 & 12/12 & 93\% \\
GPT-4-Turbo & 11/12 & 11/12 & 91\% \\
\bottomrule
\end{tabularx}\\[3pt]
\raggedright\scriptsize GPT-4-Turbo failure: GRI-Mech ignition delay. $\Agent_P$ generated a \texttt{spec.md} with explicit Euler (stiffness ratio $10^{10}$); $\Agent_J$ correctly rejected (S3 violated: explicit Euler does not converge for stiff systems) and triggered redesign, but GPT-4-Turbo's redesign selected LSODA with insufficient absolute tolerance, producing 18\% error (above the 10\% threshold). The S1--S4 verification caught the initial violation; the redesign failure is an LLM capability issue, not a verification issue.
\end{table}

\textbf{Fragility disclosure.} We cannot guarantee that future LLM updates will preserve current performance. The $\Agent_J$ S1--S4 verification gates provide a safety net---they will reject \texttt{spec.md} files that violate $\Sclass$ conditions regardless of which LLM generated them---but the pipeline's \emph{success rate} depends on $\Agent_P$'s ability to generate correct \texttt{spec.md} files within 3 redesign attempts. We recommend pinning model versions for production use and re-validating after any model update.

% =====================================================================
% EXTENDED COMPARISON AND ABLATION ON 72 PROSPECTIVE TASKS
% =====================================================================

\begin{table}[htbp]
\centering
\caption{Table S12. Extended controlled comparison on 72 prospective tasks. Without $\Sclass$-verification (No $\Agent_J$), success drops to 53\%.}
\small
\begin{tabularx}{\textwidth}{Xcccc}
\toprule
\textbf{Difficulty tier} & \textbf{Full pipeline} & \textbf{No $\Agent_J$} & \textbf{GPT-4+spec} & \textbf{PINN} \\
\midrule
Textbook (interior $\Sclass$, 18) & 18/18 (100\%) & 14/18 (78\%) & 13/18 (72\%) & 11/18 (61\%) \\
Standard (interior $\Sclass$, 30) & 27/30 (90\%) & 17/30 (57\%) & 13/30 (43\%) & 7/30 (23\%) \\
Frontier (near $\partial \Sclass$, 24) & 19/24 (79\%) & 7/24 (29\%) & 5/24 (21\%) & 1/24 (4\%) \\
\midrule
\textbf{Total} & \textbf{64/72 (89\%)} & \textbf{38/72 (53\%)} & \textbf{31/72 (43\%)} & \textbf{19/72 (26\%)} \\
\bottomrule
\end{tabularx}\\[3pt]
\raggedright\scriptsize The Judge's value increases monotonically with proximity to $\partial \Sclass$: 22 percentage-point gap on textbook problems vs.\ 50 percentage-point gap on frontier problems.
\end{table}

\begin{table}[htbp]
\centering
\caption{Table S13. Per-gate ablation of Judge Agent on 72 prospective tasks. Gates mapped to $\Sclass$ conditions S1--S4 ($n = 3$ trials each).}
\small
\begin{tabularx}{\textwidth}{Xccc}
\toprule
\textbf{Configuration} & \textbf{Correct + Bounded} & \textbf{Quality Pass} & \textbf{$\Sclass$ cond.} \\
\midrule
Full $\Agent_J$ (all gates + audit) & 64/72 (89\%) & 58/72 (81\%) & S1--S4 \\
Remove well-posedness gate & 50/72 (69\%) & 44/72 (61\%) & S2+S3 \\
Remove BC/IC compatibility & 56/72 (78\%) & 48/72 (67\%) & S2 \\
Remove dimensional consistency & 60/72 (83\%) & 52/72 (72\%) & S1 \\
Remove classification & 62/72 (86\%) & 56/72 (78\%) & S1 \\
Remove cost estimation & 64/72 (89\%) & 54/72 (75\%) & S4 feasibility \\
Remove quality audit only & 64/72 (89\%) & 52/72 (72\%) & S4 post-hoc \\
Remove all (no $\Agent_J$) & 38/72 (53\%) & 28/72 (39\%) & None \\
\bottomrule
\end{tabularx}\\[3pt]
\raggedright\scriptsize The well-posedness gate (S2+S3) remains the most critical single gate. The quality audit's contribution (S4) is most visible in the gap between ``Correct + Bounded'' and ``Quality Pass'': removing the audit does not change the bounded-error count but allows 6 additional qualitative failures (near $\partial \Sclass$, where S4 is fragile) to pass undetected.
\end{table}

% =====================================================================
% SUPPLEMENTARY FIGURES
% =====================================================================

\section*{Supplementary Figures}

\textbf{Figure S1. CT per-case PSNR distribution.} Histogram of PSNR across 200 LoDoPaB-CT test cases for the framework (blue) and expert baseline (orange). The framework distribution is slightly left-shifted (mean 31.7 vs.\ 32.1\,dB) with comparable spread ($\sigma = 1.2$ vs.\ 1.1\,dB). The 95\% CI for the mean difference is $[-0.6, -0.2]$\,dB (bootstrap, $n = 10{,}000$). Cohen's $d = 0.35$. CT reconstruction lies well within $\Sclass$ (low $L_{\text{DAG}} \approx 2.1$).

\textbf{Figure S2. Prospective benchmark quality ratio distribution.} Boxplot of quality ratio (framework / expert baseline) for the 58 fully successful prospective tasks. Median: 94\%. IQR: 89--98\%. Whiskers: 78--100\%. One outlier at 72\% (a highly optimized turbulence simulation where the expert used a custom adaptive scheme---a frontier problem near $\partial \Sclass$). Stratified by domain: imaging tasks (median 97\%), PDE tasks (93\%), optimization tasks (91\%), stochastic tasks (88\%).

\textbf{Figure S3. Scalability, $\rho$, and $\Sclass$-verification overhead.} Left panel: wall-clock time breakdown ($\Agent_P$, $\Agent_J$ $\Sclass$-verification, $\Agent_E$) across the 12 development problems, ordered by total DOF. $\Agent_J$ overhead: 3\% (heat equation, $10^3$ DOF) to 28\% (GRI-Mech, stiffness analysis dominates). $\Agent_E$ dominates above $10^4$ DOF. Right panel: $\Agent_J$ time vs.\ total time as a function of problem DOF; the ratio decreases as $O(1/\text{DOF})$ since $\Sclass$-verification cost is nearly constant while $\Agent_E$ scales with problem size. The setup-time efficiency $\rho$ is independent of $\Agent_J$ overhead, since $\rho$ measures total setup time (minutes) vs.\ expert setup time (days).

% =====================================================================
% S8: FORMAL SPEC.MD GRAMMAR AND EXTERNAL TOOL ADAPTERS
% =====================================================================

\section{Formal Grammar for \texttt{spec.md} and External Tool Adapters}
\label{sec:s8}

\subsection*{S8.1 Formal grammar}

A valid \texttt{spec.md} file is a structured Markdown document whose abstract syntax is defined by the following context-free grammar (in EBNF notation):

\begin{verbatim}
spec_file   = header , domain , equations , boundary , initial ,
              observables , tolerance , { extension } ;

header      = "# Specification:" , TITLE , NEWLINE ;

domain      = "## Domain" , NEWLINE , { key_value } ;
equations   = "## Equations" , NEWLINE , { equation_block | key_value } ;
boundary    = "## Boundary Conditions" , NEWLINE , { key_value } ;
initial     = "## Initial Conditions" , NEWLINE , { key_value } ;
observables = "## Observables" , NEWLINE , { key_value } ;
tolerance   = "## Tolerance" , NEWLINE , { key_value } ;
extension   = "## " , SECTION_NAME , NEWLINE , { key_value } ;

key_value   = KEY , ":" , VALUE , NEWLINE ;
equation_block = KEY , ":" , "|" , NEWLINE , { INDENTED_LINE } ;

KEY         = /[a-zA-Z_][a-zA-Z0-9_]*/ ;
VALUE       = /[^\n]+/ ;
TITLE       = /[^\n]+/ ;
\end{verbatim}

The six mandatory sections map one-to-one to the S1 six-tuple:
\begin{itemize}[nosep,leftmargin=1.5em]
\item \texttt{\#\# Domain} $\to \Domain$ (computational domain)
\item \texttt{\#\# Equations} $\to \mathcal{E}$ (governing equations)
\item \texttt{\#\# Boundary Conditions} $\to \mathcal{B}$ (boundary conditions)
\item \texttt{\#\# Initial Conditions} $\to \mathcal{I}$ (initial conditions)
\item \texttt{\#\# Observables} $\to \mathcal{O}$ (output functional)
\item \texttt{\#\# Tolerance} $\to \varepsilon$ (target tolerance)
\end{itemize}

Additional sections (e.g., \texttt{\#\# Primitives Required}, \texttt{\#\# Task Instance Variations}) are permitted as extensions but are not required for $\Sclass$ membership. Key-value pairs within each section use YAML-compatible syntax for machine parseability while remaining human-readable.

A \texttt{spec.md} file is \emph{valid} if and only if: (1) all six mandatory sections are present with at least one non-empty key-value pair; (2) the \texttt{Tolerance} section contains at least one numerical threshold; (3) the \texttt{Equations} section contains at least one mathematical expression (in plain text, LaTeX, or symbolic form). Validation is implemented as a 12-line Python function in the pipeline (see \texttt{validate\_specmd.py} in the repository).

\subsection*{S8.2 External tool adapters (preliminary)}

To demonstrate that \texttt{spec.md} is not locked to our pipeline, we have developed preliminary adapters that translate a subset of \texttt{spec.md} files into input formats for two external tools:

\textbf{FEniCS adapter.} For \texttt{spec.md} files describing linear elliptic PDEs (Poisson, heat equation, linear elasticity), a Python script parses the six-tuple and generates a FEniCS variational form (\texttt{.py} file). Tested on Poisson ($\nabla^2 u = f$ on $[0,1]^2$) and heat equation ($\partial_t u = \kappa \nabla^2 u$); both produce solutions matching the pipeline output to machine precision. The adapter handles: domain geometry (rectangular, from the Domain section), Dirichlet/Neumann BCs (from Boundary Conditions), initial conditions (from Initial Conditions), and target tolerance (from Tolerance, mapped to mesh refinement). It does not yet handle non-rectangular domains, nonlinear equations, or coupled systems.

\textbf{COMSOL adapter.} For the same linear elliptic subset, a script generates COMSOL \texttt{.mph} parameter files via the COMSOL--MATLAB LiveLink API. Tested on Poisson only; the generated COMSOL model produces results within $10^{-6}$ of the pipeline output. This adapter is more limited than the FEniCS adapter and requires a COMSOL license.

These adapters are proof-of-concept demonstrations covering $<5\%$ of the problems in $\Sclass$. Extending them to nonlinear, time-dependent, and coupled problems is a significant engineering effort. We include them to make the universality aspiration concrete, not to claim current universality.

% =====================================================================
\end{document}
